\scanpage{40}
% archive sheet number: 39
On en conclut un isom canonique
\[
\boxed{H^1(X,A)\mathrel{\overset{\sim}{\longleftarrow}}\mathrm{Ext}^{1}(A',\underline{\mathrm{Pic}}_{X})}
\]
\margin{Cet isomorphisme reste-t-il valable pour $X$ singulier ??? Il est possible qu'on ait besoin de $H^i(X,G_m)=0$ pour $i\geqslant2$, qui \unsure{est douteux si|est douteux en} $X$ non rég ??}
($A$, variété abélienne), qu'on comparera à la formule de dualité
\[
\boxed{H^1(X,G)\simeq\mathrm{Hom}(D(G),\underline{\mathrm{Pic}}_{X})}
\]
pour $G$ commutatif affine.

Remarque, On a des hom canoniques
\diagram{\begin{tikzcd}[ampersand replacement=\&]
\mathrm{Ext}^{1}(\underline{\mathrm{Alb}}_{X},A) \arrow[r] \arrow[d,"\wr"] \& H^1(X,A)\\
\mathrm{Ext}^{1}(A',\underline{\mathrm{Alb}}_{X}{}') \arrow[d,"\unsure{\wr|\sim}"] \& \\
\mathrm{Ext}^{1}(A',\underline{\mathrm{Pic}}^{\tau}) \&
\end{tikzcd}}
dont on montre que c'est des \uline{isomorphismes dans}, qui définissent sur les sous groupes en question l'isomorphisme envisagé plus haut.
